% !PRINTED 650
\weberpage{650}
%<<CONT
determine a functional $\psi = \varepsilon\varphi$ that is free of the $x, y, \ldots$ and
associated with $\varphi$. Since $\varphi/\psi$ is then integral, the coefficients of
$\varphi$ must, by 7., all be divisible by $\psi$, and consequently all be divisible
by $\varphi$ as well. Which was to be proved.

At the same time this proves:

9. \gesperrt{A holomorphic integral functional passes over into
an associated one when the variables are denoted
differently.}

10. \gesperrt{Every functional $\bar\omega$ in $\bar\Omega$ passes over, by multiplication
 with a unit in $\bar Z$, into a holomorphic
functional.}

For one can first represent $\bar\omega$ as the quotient of two integral
holomorphic functionals in the form
\[
\omega = \frac{\psi}{\varphi}
% [sic] printed \omega = \frac{\psi}{\varphi} without the bar over \omega, the sentence represents \bar\omega and (6) continues with \bar\omega
\]
If one puts $N(\varphi) = \varphi\varphi'$, then, as formula (4)
shows, in which $a_m E_m = \pm N(\varphi)$, $\varphi'$ is associated with a
holomorphic functional, and it follows:
\weq{6}{E_m \bar\omega = \frac{\psi \varphi'}{a_m}.}

11. \gesperrt{To an integral functional $\bar\omega$ one can choose a
second one $\mu$ in such a way that $\mu$ does not become divisible by
arbitrarily given prime functionals $\pi_1, \pi_2, \ldots$, and that}
\weq{7}{\bar\omega \mu = \alpha}
\gesperrt{becomes a function in $\Omega$.}

Let $\bar\omega$ be represented, by 10., by a holomorphic functional $\varphi$.
By 8., its coefficients are all divisible by $\bar\omega$, but not
all by $\bar\omega \pi_1$. There is therefore a function $\alpha_1$ which is divisible
by $\bar\omega$, but not by $\bar\omega \pi_1$.

Now let one put:
\[
\bar\omega_1 = \bar\omega \pi_2 \pi_3 \ldots, \quad \bar\omega_2 = \bar\omega \pi_1 \pi_3 \ldots, \quad \bar\omega_3 = \omega \pi_2 \pi_3 \ldots
% [sic] printed \bar\omega_3 = \omega \pi_2 \pi_3 \ldots, the bar over \omega is missing and the factors repeat those of \bar\omega_1; the construction requires \bar\omega_3 = \bar\omega \pi_1 \pi_2 \pi_4 \ldots
\]
and determine:
\[
\begin{aligned}
&\alpha_1 \text{ divisible by } \bar\omega_1, \text{ not divisible by } \bar\omega_1 \pi_1, \\
&\alpha_2 \text{ \textquotedbl} \text{ \textquotedbl} \bar\omega_2, \text{ \textquotedbl} \text{ \textquotedbl} \text{ \textquotedbl} \bar\omega_2 \pi_2, \\
&\alpha_3 \text{ \textquotedbl} \text{ \textquotedbl} \bar\omega_3, \text{ \textquotedbl} \text{ \textquotedbl} \text{ \textquotedbl} \bar\omega_3 \pi_3.
\end{aligned}
\]
\[
\cdots\cdots\cdots\cdots\cdots\cdots\cdots\cdots\cdots\cdots\cdots\cdots\cdots\cdots\cdots\cdots\cdots
\]
Then
\[
\alpha = \alpha_1 + \alpha_2 + \alpha_3 + \cdots
\]
satisfies the requirements of theorem 11.
